% Clase 6 --- Superficie cerrada ARBITRARIA, con y sin la carga adentro
% (§2.5 y §2.6). En (a) la superficie se descompone en cascarones cónicos como
% el de §2.3 y todos los pares se cancelan. En (b) se le resta una esferita S'
% centrada en Q: lo que queda no contiene la carga, así que Phi_S = Phi_{S'}.
\begin{center}
\begin{tikzpicture}[scale=1]

  % =============== (a) la carga queda AFUERA ===============
  \begin{scope}
    \node[figpos, minimum size=8pt] at (0,0) {};
    \node[figlbl, below left=1pt] at (0,0) {$Q$};

    % superficie arbitraria, lejos de la carga
    \draw[figline, fill=figblue!6]
      plot[smooth cycle, tension=0.8]
      coordinates {(1.8,0.95) (2.9,1.35) (3.7,0.35) (3.35,-0.95) (2.2,-1.25) (1.55,-0.3)};
    \node[figlbl, figblue, above=2pt] at (2.7,1.45) {$S$};

    % rayos desde la carga que la atraviesan
    \foreach \a in {-16,-4,8,20} {
      \draw[figvecaux] (\a:1.35) -- (\a:4.0);
    }
    \node[figlblaux, below=2pt] at (2.6,-1.6) {cada rayo entra y sale};

    \node[figlbl, figred] at (2.6,-2.45) {$\Phi = 0$};
    \node[figlbl] at (2.0,-3.15) {(a) la carga queda \textbf{afuera}};
  \end{scope}

  % =============== (b) la carga queda ADENTRO ===============
  \begin{scope}[xshift=8.0cm]
    % superficie arbitraria que la encierra
    \draw[figline, fill=figblue!6]
      plot[smooth cycle, tension=0.8]
      coordinates {(-0.3,1.7) (1.5,1.5) (2.0,0.1) (1.1,-1.5) (-0.9,-1.4) (-1.7,0.2)};
    \node[figlbl, figblue, above=2pt] at (0.6,1.75) {$S$};

    % esferita auxiliar centrada en la carga
    \draw[figred, dashed, line width=0.9pt] (0.1,0.05) circle (0.72);
    \node[figlbl, figred, below right=0pt] at (0.6,-0.45) {$S'$};

    \node[figpos, minimum size=8pt] at (0.1,0.05) {};
    \node[figlbl, above left=1pt] at (0.1,0.05) {$Q$};

    % rayos que salen
    \foreach \a in {35,105,175,250,315} {
      \draw[figvecaux] ($(0.1,0.05)+(\a:0.72)$) -- ($(0.1,0.05)+(\a:2.15)$);
    }

    \node[figlbl, figred, align=center] at (0.2,-2.4)
      {$\Phi_S = \Phi_{S'} = \dfrac{Q}{\varepsilon_0}$};
    \node[figlbl] at (0.2,-3.15) {(b) la carga queda \textbf{adentro}};
  \end{scope}

\end{tikzpicture}\\[0.5em]
{\small\itshape En (a) la superficie se descompone en cascarones cónicos como
el de la figura anterior: cada trozo de entrada cancela al de salida, sin
importar la forma. En (b) el truco es \textbf{restar} una esferita $S'$ centrada
en la carga: la región entre $S$ y $S'$ es cerrada y \textbf{no} contiene a $Q$,
así que su flujo es nulo y por lo tanto $S$ tiene el mismo flujo que $S'$, que
ya sabemos calcular.}
\end{center}
